\documentclass[11pt,a4paper]{article}

% Packages
\usepackage[utf8]{inputenc}
\usepackage[T1]{fontenc}
\usepackage{geometry}
\usepackage{graphicx}
\usepackage{hyperref}
\usepackage{booktabs}
\usepackage{enumitem}
\usepackage{float}
\usepackage{fancyhdr}
\usepackage{parskip}
\usepackage{titlesec}
\usepackage{listings}
\usepackage{xcolor}
\usepackage{longtable}

% Page geometry
\geometry{left=2.5cm, right=2.5cm, top=2.5cm, bottom=2.5cm}

% Code listing style
\lstdefinestyle{json}{
    basicstyle=\ttfamily\small,
    breaklines=true,
    frame=single,
    backgroundcolor=\color{gray!10},
    keywordstyle=\color{blue},
    stringstyle=\color{red!70!black},
    showstringspaces=false,
}

\lstdefinestyle{bash}{
    basicstyle=\ttfamily\small,
    breaklines=true,
    frame=single,
    backgroundcolor=\color{black!5},
}

% Header
\pagestyle{fancy}
\fancyhf{}
\rhead{Smart Productivity Analytics Platform}
\lhead{API Documentation}
\rfoot{Page \thepage}

% Hyperlink setup
\hypersetup{colorlinks=true, linkcolor=blue, urlcolor=blue}

\title{\textbf{Smart Productivity Analytics Platform}\\[0.3cm]\Large API Documentation\\[0.2cm]\normalsize Version 1.0.0}
\author{COMP3011 Web Services and Web Data}
\date{March 2026}

\begin{document}

\maketitle
\thispagestyle{fancy}

\tableofcontents
\newpage

% =============================================================================
\section{Introduction}
% =============================================================================

This document provides comprehensive documentation for the Smart Productivity Analytics Platform API. The API enables task management and productivity analytics through a RESTful interface.

\subsection{Base URL}

\begin{lstlisting}[style=bash]
https://your-domain.com/api/v1/
\end{lstlisting}

For local development:
\begin{lstlisting}[style=bash]
http://localhost:8000/api/v1/
\end{lstlisting}

\subsection{API Versioning}

The API uses URL-based versioning. All endpoints are prefixed with \texttt{/api/v1/}.

\subsection{Content Type}

All requests and responses use JSON format:
\begin{lstlisting}[style=bash]
Content-Type: application/json
\end{lstlisting}

% =============================================================================
\section{Authentication}
% =============================================================================

The API uses JWT (JSON Web Token) authentication. Tokens must be included in the \texttt{Authorization} header.

\subsection{Token Format}

\begin{lstlisting}[style=bash]
Authorization: Bearer <access_token>
\end{lstlisting}

\subsection{Token Lifecycle}

\begin{table}[H]
\centering
\begin{tabular}{@{}lll@{}}
\toprule
\textbf{Token Type} & \textbf{Lifetime} & \textbf{Purpose} \\
\midrule
Access Token & 60 minutes & API request authentication \\
Refresh Token & 7 days & Obtain new access tokens \\
\bottomrule
\end{tabular}
\caption{JWT Token Configuration}
\end{table}

\subsection{Obtaining Tokens}

\textbf{Endpoint:} \texttt{POST /api/v1/auth/token/}

\textbf{Request:}
\begin{lstlisting}[style=json]
{
    "email": "user@example.com",
    "password": "your_password"
}
\end{lstlisting}

\textbf{Response (200 OK):}
\begin{lstlisting}[style=json]
{
    "access": "eyJ0eXAiOiJKV1QiLCJhbGc...",
    "refresh": "eyJ0eXAiOiJKV1QiLCJhbGc..."
}
\end{lstlisting}

\subsection{Refreshing Tokens}

\textbf{Endpoint:} \texttt{POST /api/v1/auth/token/refresh/}

\textbf{Request:}
\begin{lstlisting}[style=json]
{
    "refresh": "<refresh_token>"
}
\end{lstlisting}

% =============================================================================
\section{Rate Limiting}
% =============================================================================

\begin{table}[H]
\centering
\begin{tabular}{@{}ll@{}}
\toprule
\textbf{User Type} & \textbf{Limit} \\
\midrule
Anonymous & 100 requests/hour \\
Authenticated & 1000 requests/hour \\
\bottomrule
\end{tabular}
\caption{Rate Limits}
\end{table}

When rate limited, the API returns HTTP 429 with a \texttt{Retry-After} header.

% =============================================================================
\section{Error Handling}
% =============================================================================

\subsection{HTTP Status Codes}

\begin{table}[H]
\centering
\begin{tabular}{@{}lll@{}}
\toprule
\textbf{Code} & \textbf{Status} & \textbf{Description} \\
\midrule
200 & OK & Request successful \\
201 & Created & Resource created successfully \\
204 & No Content & Resource deleted successfully \\
400 & Bad Request & Invalid request data \\
401 & Unauthorized & Authentication required \\
403 & Forbidden & Permission denied \\
404 & Not Found & Resource not found \\
429 & Too Many Requests & Rate limit exceeded \\
500 & Internal Server Error & Server error \\
\bottomrule
\end{tabular}
\caption{HTTP Status Codes}
\end{table}

\subsection{Error Response Format}

\begin{lstlisting}[style=json]
{
    "detail": "Error message describing the issue",
    "code": "error_code"
}
\end{lstlisting}

% =============================================================================
\section{Endpoints Reference}
% =============================================================================

% -----------------------------------------------------------------------------
\subsection{Authentication Endpoints}
% -----------------------------------------------------------------------------

\begin{table}[H]
\centering
\small
\begin{tabular}{@{}llp{6cm}@{}}
\toprule
\textbf{Method} & \textbf{Endpoint} & \textbf{Description} \\
\midrule
POST & /auth/register/ & Register a new user account \\
POST & /auth/token/ & Obtain JWT token pair \\
POST & /auth/token/refresh/ & Refresh access token \\
POST & /auth/token/blacklist/ & Invalidate refresh token (logout) \\
GET & /auth/profile/ & Get current user profile \\
PUT/PATCH & /auth/profile/ & Update user profile \\
POST & /auth/change-password/ & Change password \\
\bottomrule
\end{tabular}
\caption{Authentication Endpoints}
\end{table}

\subsubsection{Register User}

\textbf{POST} \texttt{/api/v1/auth/register/}

\textbf{Request Body:}
\begin{lstlisting}[style=json]
{
    "email": "user@example.com",
    "username": "johndoe",
    "password": "SecurePassword123!",
    "password_confirm": "SecurePassword123!"
}
\end{lstlisting}

\textbf{Response (201 Created):}
\begin{lstlisting}[style=json]
{
    "user": {
        "id": "550e8400-e29b-41d4-a716-446655440000",
        "email": "user@example.com",
        "username": "johndoe"
    },
    "tokens": {
        "access": "eyJ0eXAiOiJKV1Q...",
        "refresh": "eyJ0eXAiOiJKV1Q..."
    },
    "message": "Registration successful."
}
\end{lstlisting}

% -----------------------------------------------------------------------------
\subsection{Task Endpoints}
% -----------------------------------------------------------------------------

\begin{table}[H]
\centering
\small
\begin{tabular}{@{}llp{6cm}@{}}
\toprule
\textbf{Method} & \textbf{Endpoint} & \textbf{Description} \\
\midrule
GET & /tasks/ & List all tasks (with filtering) \\
POST & /tasks/ & Create a new task \\
GET & /tasks/\{id\}/ & Get task details \\
PUT & /tasks/\{id\}/ & Update task (full) \\
PATCH & /tasks/\{id\}/ & Update task (partial) \\
DELETE & /tasks/\{id\}/ & Delete task \\
POST & /tasks/\{id\}/complete/ & Mark task as completed \\
POST & /tasks/\{id\}/reopen/ & Reopen completed task \\
GET & /tasks/today/ & Get today's tasks \\
GET & /tasks/overdue/ & Get overdue tasks \\
POST & /tasks/bulk\_update\_status/ & Bulk update task status \\
\bottomrule
\end{tabular}
\caption{Task Endpoints}
\end{table}

\subsubsection{Task Object}

\begin{lstlisting}[style=json]
{
    "id": "550e8400-e29b-41d4-a716-446655440000",
    "title": "Complete project report",
    "description": "Write the final section",
    "status": "pending",
    "priority": 4,
    "energy_level": "high",
    "due_date": "2026-03-15T17:00:00Z",
    "category": {
        "id": "...",
        "name": "Work"
    },
    "tags": [
        {"id": "...", "name": "urgent"}
    ],
    "created_at": "2026-03-10T10:00:00Z",
    "updated_at": "2026-03-10T10:00:00Z"
}
\end{lstlisting}

\subsubsection{Create Task}

\textbf{POST} \texttt{/api/v1/tasks/}

\textbf{Request Body:}
\begin{lstlisting}[style=json]
{
    "title": "Complete project report",
    "description": "Write the final section",
    "priority": 4,
    "energy_level": "high",
    "due_date": "2026-03-15T17:00:00Z",
    "category": "550e8400-e29b-41d4-a716-446655440000",
    "tags": ["550e8400-e29b-41d4-a716-446655440001"]
}
\end{lstlisting}

\subsubsection{Query Parameters for List}

\begin{table}[H]
\centering
\small
\begin{tabular}{@{}llp{5cm}@{}}
\toprule
\textbf{Parameter} & \textbf{Type} & \textbf{Description} \\
\midrule
status & string & Filter by status (pending, in\_progress, completed, cancelled) \\
priority & integer & Filter by priority (1-5) \\
category & UUID & Filter by category ID \\
energy\_level & string & Filter by energy level (low, medium, high) \\
search & string & Search in title and description \\
ordering & string & Sort field (e.g., -due\_date, priority) \\
\bottomrule
\end{tabular}
\caption{Task Query Parameters}
\end{table}

% -----------------------------------------------------------------------------
\subsection{Category Endpoints}
% -----------------------------------------------------------------------------

\begin{table}[H]
\centering
\small
\begin{tabular}{@{}llp{6cm}@{}}
\toprule
\textbf{Method} & \textbf{Endpoint} & \textbf{Description} \\
\midrule
GET & /categories/ & List all categories \\
POST & /categories/ & Create a category \\
GET & /categories/\{id\}/ & Get category details \\
PUT/PATCH & /categories/\{id\}/ & Update category \\
DELETE & /categories/\{id\}/ & Delete category \\
GET & /categories/\{id\}/stats/ & Get category statistics \\
\bottomrule
\end{tabular}
\caption{Category Endpoints}
\end{table}

% -----------------------------------------------------------------------------
\subsection{Tag Endpoints}
% -----------------------------------------------------------------------------

\begin{table}[H]
\centering
\small
\begin{tabular}{@{}llp{6cm}@{}}
\toprule
\textbf{Method} & \textbf{Endpoint} & \textbf{Description} \\
\midrule
GET & /tags/ & List all tags \\
POST & /tags/ & Create a tag \\
GET & /tags/\{id\}/ & Get tag details \\
PUT/PATCH & /tags/\{id\}/ & Update tag \\
DELETE & /tags/\{id\}/ & Delete tag \\
\bottomrule
\end{tabular}
\caption{Tag Endpoints}
\end{table}



% -----------------------------------------------------------------------------
\subsection{Analytics Endpoints}
% -----------------------------------------------------------------------------

\begin{table}[H]
\centering
\small
\begin{tabular}{@{}llp{6cm}@{}}
\toprule
\textbf{Method} & \textbf{Endpoint} & \textbf{Description} \\
\midrule
GET & /analytics/dashboard/ & Dashboard overview \\
GET & /analytics/tasks/ & Task analytics \\
GET & /analytics/trends/ & Productivity trends \\
\bottomrule
\end{tabular}
\caption{Analytics Endpoints}
\end{table}

\subsubsection{Dashboard Response}

\begin{lstlisting}[style=json]
{
    "today": {
        "tasks_created": 5,
        "tasks_completed": 3
    },
    "week": {
        "tasks_created": 25,
        "tasks_completed": 18,
        "completion_rate": 72.0
    },
    "recent_tasks": [
        {
            "id": "...",
            "title": "Write report",
            "status": "completed"
        }
    ]
}
\end{lstlisting}

% -----------------------------------------------------------------------------
\subsection{MCP (Model Context Protocol) Endpoints}
% -----------------------------------------------------------------------------

MCP endpoints enable AI agent integration with standardized tool discovery and execution.

\begin{table}[H]
\centering
\small
\begin{tabular}{@{}llp{6cm}@{}}
\toprule
\textbf{Method} & \textbf{Endpoint} & \textbf{Description} \\
\midrule
GET & /mcp/capabilities/ & Get API capabilities (public) \\
GET & /mcp/tools/ & List available tools \\
POST & /mcp/execute/ & Execute a tool \\
\bottomrule
\end{tabular}
\caption{MCP Endpoints}
\end{table}

\subsubsection{Available Tools}

\begin{table}[H]
\centering
\small
\begin{tabular}{@{}lp{8cm}@{}}
\toprule
\textbf{Tool Name} & \textbf{Description} \\
\midrule
create\_task & Create a new task \\
list\_tasks & List tasks with filtering \\
complete\_task & Mark a task as completed \\
get\_today\_tasks & Get tasks due today \\
get\_overdue\_tasks & Get overdue tasks \\
suggest\_next\_task & Get AI-powered task suggestion based on energy level \\
get\_productivity\_summary & Get productivity metrics \\
create\_category & Create a new category \\
\bottomrule
\end{tabular}
\caption{MCP Tools}
\end{table}

\subsubsection{Execute Tool}

\textbf{POST} \texttt{/api/v1/mcp/execute/}

\textbf{Request Body:}
\begin{lstlisting}[style=json]
{
    "tool": "suggest_next_task",
    "arguments": {
        "energy_level": "low"
    }
}
\end{lstlisting}

\textbf{Response (200 OK):}
\begin{lstlisting}[style=json]
{
    "success": true,
    "tool": "suggest_next_task",
    "result": {
        "task": {
            "id": "...",
            "title": "Organize documents",
            "priority": 2
        },
        "reason": "Low-energy task suitable for current state"
    }
}
\end{lstlisting}

% =============================================================================
\section{Data Models}
% =============================================================================

\subsection{Priority Levels}

\begin{table}[H]
\centering
\begin{tabular}{@{}ll@{}}
\toprule
\textbf{Value} & \textbf{Label} \\
\midrule
1 & Lowest \\
2 & Low \\
3 & Medium \\
4 & High \\
5 & Urgent \\
\bottomrule
\end{tabular}
\caption{Task Priority Levels}
\end{table}

\subsection{Task Status}

\begin{table}[H]
\centering
\begin{tabular}{@{}ll@{}}
\toprule
\textbf{Value} & \textbf{Description} \\
\midrule
pending & Task not yet started \\
in\_progress & Task is being worked on \\
completed & Task is finished \\
cancelled & Task was cancelled \\
\bottomrule
\end{tabular}
\caption{Task Status Values}
\end{table}

\subsection{Energy Levels}

\begin{table}[H]
\centering
\begin{tabular}{@{}ll@{}}
\toprule
\textbf{Value} & \textbf{Suitable For} \\
\midrule
low & Administrative tasks, organizing \\
medium & Standard work tasks \\
high & Complex problem-solving, creative work \\
\bottomrule
\end{tabular}
\caption{Energy Level Values}
\end{table}



% =============================================================================
\section{Examples}
% =============================================================================

\subsection{Complete Workflow Example}

\subsubsection{1. Register and Authenticate}

\begin{lstlisting}[style=bash]
# Register
curl -X POST http://localhost:8000/api/v1/auth/register/ \
  -H "Content-Type: application/json" \
  -d '{"email":"user@example.com","username":"john",
       "password":"SecurePass123!","password_confirm":"SecurePass123!"}'

# Login
curl -X POST http://localhost:8000/api/v1/auth/token/ \
  -H "Content-Type: application/json" \
  -d '{"email":"user@example.com","password":"SecurePass123!"}'
\end{lstlisting}

\subsubsection{2. Create and Manage Tasks}

\begin{lstlisting}[style=bash]
# Create task
curl -X POST http://localhost:8000/api/v1/tasks/ \
  -H "Authorization: Bearer <token>" \
  -H "Content-Type: application/json" \
  -d '{"title":"Write report","priority":4,"due_date":"2026-03-15T17:00:00Z"}'

# Complete task
curl -X POST http://localhost:8000/api/v1/tasks/<id>/complete/ \
  -H "Authorization: Bearer <token>"
\end{lstlisting}

\subsubsection{3. Check Analytics}

\begin{lstlisting}[style=bash]
# Get dashboard
curl -X GET http://localhost:8000/api/v1/analytics/dashboard/ \
  -H "Authorization: Bearer <token>"

# Get productivity trends
curl -X GET http://localhost:8000/api/v1/analytics/trends/?period=week \
  -H "Authorization: Bearer <token>"
\end{lstlisting}

% =============================================================================
\section{OpenAPI Specification}
% =============================================================================

The complete OpenAPI 3.0 specification is available at:

\begin{itemize}
    \item \textbf{Interactive (Swagger UI):} \texttt{/api/v1/docs/swagger/}
    \item \textbf{ReDoc:} \texttt{/api/v1/docs/redoc/}
    \item \textbf{JSON Schema:} \texttt{/api/v1/docs/schema/}
    \item \textbf{Repository:} \url{https://github.com/Ejkenway/COMP3011-CW1/blob/main/todo_api/openapi.yaml}
\end{itemize}

\end{document}
